\section{相关研究现状}

\subsection{基于工具调用的大语言模型智能体（LLM Agent）研究现状}

\subsubsection{面向工具调用的大语言模型技术演进脉络}

从技术演进脉络看，面向智能体的大语言模型发展可分为预训练（Pre-training）、适配调优（Adaptation Tuning）与利用（Utilization）三大技术维度\cite{zhao2023survey}，其演进历程可概括为三个关键阶段：\textbf{第一阶段（2017--2019）}是架构奠基与预训练范式确立期，Vaswani等\cite{vaswani2017attention}提出Transformer架构，以自注意力机制取代循环结构，奠定了现代语言模型的计算基础；Radford等\cite{radford2018improving}与Devlin等\cite{devlin2019bert}分别通过GPT-1和BERT验证了预训练-微调范式在语言理解任务上的有效性，确立了语言模型“大规模预训练+任务微调”的基本范式。\textbf{第二阶段（2020--2022）}是规模扩张与能力对齐期，OpenAI\cite{brown2020language}推出1750亿参数的GPT-3，首次展示了大语言模型在少样本（Few-shot）场景下的强大泛化能力；Kaplan等\cite{kaplan2020scaling}系统揭示了模型性能与参数量、数据量、计算量之间的幂律缩放关系（Scaling Laws），为后续模型规模扩张提供了可预测的理论依据；Wei等\cite{wei2022emergent}发现当模型参数规模超过特定阈值后，会涌现出小模型不具备的能力，如多步推理、指令遵循和上下文学习，随后据此提出了思维链（Chain-of-Thought，COT）\cite{bosma2022chainofthought}技术，通过在提示词中展示中间推理步骤来激发大型语言模型推理能力，证明了足够大的语言模型在经过少样本CoT提示后，无需微调即可在算术、常识和符号推理任务上取得大幅性能提升；Ouyang等\cite{ouyang2022training}提出InstructGPT，通过指令微调与基于人类反馈的强化学习使模型行为与人类意图对齐，解决了大规模模型“能力强但不可控”的核心矛盾，标志着LLM从能力涌现走向能力可控的关键转折；Schick等\cite{schick2023toolformer}提出的Toolformer探索了模型自监督学习何时调用API、调用何种API以及如何吸收工具返回结果，证明仅需少量标注样本即可使语言模型掌握多种外部工具的使用，从而大幅降低了工具调用能力的训练成本；Wei等\cite{wei2022chain}提出思维链（Chain-of-Thought）提示方法，通过引导模型生成中间推理步骤显著提升多步推理能力，为LLM执行复杂规划任务奠定了推理基础。\textbf{第三阶段（2023至今）}是模型层面的Agent能力集成期，人们对大语言模型的利用维度从文本生成逐渐扩展至Agent式任务执行。OpenAI发布的GPT-4\cite{openai2023gpt4}不仅在多项专业和学术基准上达到人类水平表现，更首次原生支持函数调用（Function Calling）机制，使LLM能够直接生成结构化的工具调用请求并解析返回结果，标志着LLM从文本生成器向任务执行器的技术转变；Gorilla\cite{patil2024gorilla}通过在大规模API语料上微调LLM并使其自主选择并调用真实API，首次验证了LLM可在不依赖人工提示模板的情况下实现高精度工具选择；Qin等\cite{qin2024toolllm}提出的ToolLLM通过构建大规模自动化指令数据集实现了复杂多工具场景的灵活应用，并引入深度优先搜索决策树策略以提升长链工具调用的探索效率与成功率，为开源大模型在真实复杂任务中的工具编排与决策提供了系统化方案；Meta推出的LLaMA\cite{touvron2023llama}和阿里巴巴推出的Qwen\cite{yang2025qwen3technicalreport}系列大模型则使工具调用能力的训练与适配可在开源模型上复现，推动了Agent研究的可复现性与民主化。至此，指令遵循与多步推理能力的涌现与对齐，使得LLM不再仅作为单纯的文本生成器，而是具备了理解任务意图、分解执行步骤并调用外部工具的潜力，这为LLM从语言模型向智能体的角色跃迁奠定了完整的能力基础。

\subsubsection{基于大语言模型的智能体（Agent）的演化历程}

在\textbf{Agent定义}层面，Wooldridge和Jennings\cite{wooldridge1995intelligent}于1995年从软件智能体角度提出Agent应具备自主性（Autonomy）、社会性（Social Ability）、反应性（Reactivity）和预动性（Pro-activeness）四项核心属性，而Russell和Norvig\cite{russell2020artificial}则于2020年在经典教材《Artificial Intelligence: A Modern Approach》中明确将智能体定义为“通过传感器感知环境并通过执行器作用于环境的实体”。在\textbf{LLM Agent领域}，Weng\cite{weng2023agent}提出了著名的Agent开发公式（Agent=大模型+规划+工具调用+记忆），奠定了现今LLM Agent的形态基础（如\autoref{fig:agent-overview}所示）；Xi等\cite{xi2025rise}将LLM Agent定义为以大语言模型为核心大脑（Brain），结合感知（Perception）与动作（Action）模块，能够在动态环境中自主感知、推理、规划并执行任务的智能体系统，并提出了Brain-Perception-Action三组件通用框架；Wang等\cite{wang2024survey}进一步指出，LLM Agent区别于传统软件智能体的关键在于其以自然语言为统一接口，通过提示词（Prompt）驱动推理与行动，而非依赖预编程规则Masterman等\cite{masterman2024landscape}归纳了AI Agent架构，总结了推理、规划与工具调用三大核心能力的设计范式:推理以思考-行动交织与自反思为主，规划包括任务分解、多计划选择与反思修正等方法，工具调用则围绕单Agent串行（规划、选择与反思）与多Agent分工（职责分离、Agent团队组织）两种模式展开。。而在Agent技术的实践落地层面，Parisi等\cite{parisi2022talm}提出了工具增强语言模型，通过纯文本到文本的API调用格式将外部工具接入语言模型，并设计迭代式“自博弈”机制使模型从少量工具演示中自举学习“何时调用工具及如何构造参数”，证明了工具增强可作为降低语言模型对规模依赖的互补路径；随后，Yao等\cite{yao2023react}提出了ReAct（Reason-Action）提示范式，将推理与行动交织结合，使模型能够交替生成推理轨迹与外部动作（如工具调用），通过与外部环境交互获取实时信息来指导推理与决策过程，有效缓解了纯思维链推理易产生事实幻觉的问题，为LLM向智能体演进提供了早期的交互范式；Shinn等\cite{shinn2023reflexion}则进一步提出了Reflexion机制，旨在通过语言反馈机制使Agent能够从失败经验中学习，将错误调用经验转化为自然语言反思并存入记忆，进一步提升了多步工具调用的鲁棒性与自适应能力；Wang等提出的\cite{wang2024codeact}CodeAct将工具调用动作统一为可执行Python代码，替代传统JSON格式调用，使单次动作即可完成多步操作与条件分支，在减少交互轮次的同时并显著提升了任务完成率。这些工作开始让大语言模型突破传统文本生成的形态限制，让LLM不再仅依赖内部参数完成推理，而是通过外部API、检索系统、文件系统、邮件日程、数据库、浏览器或代码执行器等工具获取信息并改变外部环境，成为与环境交互甚至嵌入环境本身的实体。在\textbf{多Agent协作}层面，Qian等\cite{qian2024chatdev}提出ChatDev，以“首席执行官（CEO）—首席技术官（CTO）—程序员—测试员”等软件公司角色链驱动多Agent对话式协作完成软件开发，验证了角色分工式对话在代码生成任务上的可行性；Wu等\cite{wu2024autogen}提出AutoGen框架，将Agent抽象为可定制、可对话的实体，支持LLM推理、人类输入和工具调用的灵活组合模式，允许开发者通过自然语言与代码混合编程定义Agent间的多轮对话拓扑、自动回复策略和人工介入条件，使得多Agent系统的编排从手工硬编码转向声明式配置，已被广泛应用于数学推理、代码生成和运营决策等场景；Hong等\cite{hong2024metagpt}提出MetaGPT框架，将软件公司的标准作业流程（Standard Operating Procedure, SOP）编码为提示序列，定义产品经理、架构师、工程师和QA工程师等多重角色，通过共享消息池与发布-订阅机制实现结构化通信，Agent间以文档和图表等结构化产出物而非自由对话进行交接，验证了SOP驱动的多Agent系统在复杂软件工程任务上的有效性。然而，上述框架在设计时均以任务完成率为核心优化目标，缺乏对工具调用链的过程级安全审计机制——ChatDev的纯对话式协作使注入内容可在角色间无约束传播，AutoGen的灵活对话拓扑意味着恶意Agent可通过多轮对话逐步提升权限，MetaGPT的结构化SOP虽降低了幻觉风险但未对跨角色数据流进行安全隔离，这些安全缺口的客观存在正是本研究所关注的问题。自动程序从“规则驱动”到“语言驱动”的范式转变，正是LLM Agent安全问题的根源所在：当智能体的行为由语义而非语法决定时，传统的访问控制与权限管理机制将面临根本性的适用性挑战\cite{chu2026systematic}。

\subsection{LLM Agent工具调用链轨迹溯源研究现状}
\todo{补充Agent行为轨迹溯源的方法和粒度；攻击向量图重绘}
% https://applink.feishu.cn/client/message/link/open?token=AmnrGOiNgYy%2Baj%2BKWJoZC8k%3D

从系统安全的视角看，将工具调用执行轨迹结构化为溯源图（Provenance Graph），是解决上述问题的天然路径。前述防御方法虽然在各自设计的攻击场景下取得了不同程度的成功，但均以单次调用或单次交互为粒度进行安全判断，缺乏对工具调用链整体结构的建模能力，即使每次工具调用均经过授权，但授权动作的语义组合可能超越预期授权范围，而这一过程对基于单次调用的访问控制系统完全不可见，这正是“语义权限提升”等跨步骤组合型威胁能够绕过现有防御的结构性原因\cite{chu2026systematic}。传统系统安全领域长期使用审计日志和溯源图刻画进程、文件、网络连接之间的因果关系\cite{gao2021threatraptor,cheng2024kairos,inam2023sok}，其核心价值在于能够表达跨进程、跨时间的因果依赖关系，从而支持从异常事件出发沿依赖边后向追溯攻击根因与影响范围。对于LLM Agent而言，提示词（Prompt）、计划（Plan）、工具调用（Tool Call）、观察（Observation）、记忆（Memory）、最终答案（Final Answer）等执行片段同样构成可追踪的行为链条，将这些片段结构化为异构图节点，并以数据流、控制流、权限依赖和时序关系为边，可以同时表达“谁触发了哪个工具”“参数来源于何处”“敏感数据经过哪些转换”“最终风险由哪些路径贡献”“防御动作在何处介入”等信息。这种图化的统一表示有望实现三个关键能力：其一，在过程层面量化组合风险，将工具敏感度、权限扩散度、调用深度和不可逆操作比例等指标嵌入图结构中进行路径级计算，弥补现有结果导向指标的不足；其二，在结构层面检测跨步骤威胁，通过图模式匹配或图异常检测识别“读取敏感文件—摘要压缩—调用外部发送工具”等链式风险路径，填补单步防御的结构性盲区；其三，在证据层面支持可解释归因，从异常节点沿依赖边后向追溯，定位风险路径的构成节点、关键触发工具和参数来源，为定向防御部署提供可信证据链\cite{wang2026fromagenttraces,dewitt2026openchallenges}。文献\cite{an2025ipiguard}中基于工具依赖图的检测方法已初步验证了这一方向的有效性，但该方法仅建模工具间的正常依赖模式，尚未拓展至包含数据流、权限状态和防御动作的完整行为轨迹图；文献\cite{he2025sentinelagent}和\cite{wang2026aligning}中的工作尝试将图表示引入Agent异常检测和授权验证，但仍局限于单Agent场景或依赖预定义策略，缺乏对多Agent委托链和跨权限域风险路径的统一建模能力。因此，如何将传统溯源图的方法论体系与LLM Agent的执行语义有机结合，构建适应Agent工具调用链特征的异构轨迹图表示，并在此基础上实现面向跨步骤组合风险的威胁检测、可解释归因与最小干预防御，是亟待解决的核心研究问题。

综合以上研究现状可以看出：LLM Agent工具调用能力已从学术原型进入大规模产业部署阶段，但现有能力评测基准以任务完成率为核心指标，较少关注工具调用链的过程合规性与安全性；攻击面已从单次提示注入扩展至工具误用、权限失配、持久化注入和供应链污染等多维威胁，且多向量组合攻击的成功率远超单点攻击；防御方法已在策略执行与系统设计、运行时检查等方向取得了阶段性进展，但大多针对特定攻击场景，缺乏通用的过程级安全审计能力，且面临安全性与任务效用之间的深层权衡；风险评估与轨迹溯源研究已初步探索了图化表示和轨迹审计的技术路线，但在跨步骤组合风险度量、可解释威胁归因和最小干预防御三个维度上仍处于起步阶段。

\subsection{LLM Agent工具调用链安全评估与防御研究现状}

针对LLM Agent技术发展带来的伴生安全性问题的研究已成为学界的一个新兴方向。工具调用能力的引入使LLM Agent的安全问题发生了本质变化。与早期基于越狱的大模型安全研究\cite{wei2023jailbreak,zou2023universal,mehrotra2024tree,ding2024wolf,chao2025jailbreaking,andriushchenko2025jailbreaking}更多关注模型是否直接生成有害信息（例如，违法犯罪、种族歧视、侮辱性言论）不同，当前基于工具调用的Agent的风险在于语言决策会被进一步转化为读取、写入、删除、发送、查询、执行命令和控制设备等外部动作。随着LLM与外部环境交互能力的增强，随之而来的安全与治理的挑战将成倍增加\cite{weidinger2022taxonomy}。多个Agent能力评测基准\cite{zhou2023webarena,deng2023mind2web,xie2024osworld,qin2024toolllm,bonatti2025windows}已显示，真实Agent任务通常具有长时序、多步骤、跨应用和跨工具组合特征，Agent执行过程通常为“任务理解—计划生成—工具调用—结果观察—状态更新—继续调用”的链式过程；工具调用的每一次执行都构成具有现实侧效应的动作，且工具返回结果通常在无显式信任标记的情况下被注入Agent规划上下文，使得数据流污染、权限误用和不可逆操作等问题从推理层延伸至执行层\cite{chu2026systematic}；LLM序列规划中的单次错误会沿推理链持续放大与传播，最终导致Agent行为偏离用户初始指令，导致“级联失效”（偏离用户任务意图）\cite{deng2025agentsunderthreat}、“语义权限提升”或“链式攻击”（Agent非法调用了用户任务中无权调用的工具）\cite{chu2026systematic,dewitt2026openchallenges}。即使Agent运行结果正确，这个结果也可能通过不必要的、不安全的或违反策略的工具调用产生，而错误的结果则可能源于误导性证据、错误的工具输出、过时的记忆、格式错误的参数、无效的中间声明、环境误观察或多智能体协调失败\cite{wang2026fromagenttraces}。在Agent链中，当多个Agent以管道形式协作时，确定哪个Agent执行了特定操作、追踪完整的委托链条在技术上极为困难，导致现有基于角色的访问控制（Role-Based Access Control, RBAC）、基于属性的访问控制（Attribute-Based Access Control, ABAC）等人类身份认证与安全管理机制在Agent场景下失效\cite{dewitt2026openchallenges}。这意味着即使单个工具调用看似符合权限要求，当多个调用按照特定顺序组合后，也可能形成敏感数据外传、权限滥用、信息污染甚至委托失控等组合型风险，Agent调用链越长、工具越多、权限越复杂，越容易出现系统级风险。同时，由于Agent的CoT轨迹本质上是基于纯文本的无状态数据，容易与LLM Agent底层的实际计算过程相脱耦，虽然CoT轨迹正越来越多地被用作Agent行为的审计日志，但若发生这种脱耦，基于CoT的审计将无法提供可靠的行为溯源与问责基础\cite{chu2026systematic}。

Greshake等\cite{greshake2023nowwhat}于2023年首次系统揭示了LLM集成应用中的间接提示注入（Indirect Prompt Injection）威胁，攻击者可通过向网页、邮件、文档等外部数据源嵌入恶意自然语言指令，使Agent在检索和处理外部内容时被诱导执行非预期操作，从而揭示了工具调用场景下“数据-指令边界模糊”这一结构性安全缺陷。此后，间接提示注入迅速成为Agent安全领域最活跃的研究议题之一：Debenedetti等\cite{balunovic2024agentdojo}提出的AgentDojo以真实世界的多轮工具调用任务为评估场景，为该方向奠定了动态基准基础；Ruan等\cite{ruan2024toolemu}在提出的ToolEmu则通过LM仿真沙箱在无副作用条件下对工具调用风险进行规模化预评估。这些早期基准的揭示了工具调用将Agent的攻击面从模型内部的生成安全性扩展至模型与外部世界交互过程中的行为安全性。

随着Agent所集成的工具数量快速增长、工具调用链的深度和组合复杂度持续增加，攻击向量迅速从内容级注入向工具全生命周期拓展。在工具选择层面，Shi等\cite{shi2026pitoolselection}提出的ToolHijacker首次将针对工具选择决策的注入攻击形式化为独立威胁面，揭示了现有安全对齐技术在Agent工具调用场景中的系统性失效边界。在工具权限层面，Zhang等\cite{zhang2026evaluatingprivilege}的实证研究指出Agent在真实世界工具使用中请求的权限远超任务必需；Jha等\cite{jha2026agent}进一步发现64.79\%的环境错误场景可触发越权工具调用，且Agent能力越强，越权行为越严重。在工具供应链层面，Tie等\cite{tie2026badskill}提出的BadSkill展示了技能投毒后门高达97.5\%--99.5\%的注入成功率，攻击者可在技能发布前将恶意行为编码为触发式后门；Wu等\cite{wu2026chainfuzzer}提出的ChainFuzzer在主流工具链中发现365个漏洞，其中82.74\%需多工具路径方可触发，揭示了工具链的组合性引入了一类难以被单步安全检查所捕获的隐式安全风险。Deng等\cite{deng2025agentsunderthreat}在ACM Computing Surveys的综述中将上述威胁归纳为Agent-to-Tool攻击维度的主动与被动二分法，明确将工具调用确立为Agent安全的独立核心威胁面；Dehghantanha等\cite{dehghantanha2026sokattacksurfaceagentic}的SoK研究进一步将工具提升为与RAG、自主性并列的三维攻击面之一，并指出多向量组合攻击的威胁等级远超单点攻击。

面对日益严峻的工具调用安全威胁，学术界与工业界在上述模型内建防御之外，进一步探索了递进的系统层防御范式，按其设计哲学可大致划分为四个层次。\textbf{第一层为模型层防御}，即通过对底层语言模型进行安全微调或对齐训练，使模型内化抵抗注入攻击的行为偏好。Wallace等\cite{wallace2024instruction}提出指令层级（Instruction Hierarchy）训练方法，使模型在遇到冲突时优先遵循系统高优先级指令而非用户或工具注入的低优先级指令，为此类防御奠定了训练范式基础；Chen等\cite{chen2025secalign}的SecAlign利用直接偏好优化（DPO）训练模型区分并拒绝注入攻击，在AgentDojo基准上平均攻击成功率降至11.77\%；Chen等\cite{chen2025metasecalign}进一步发布Meta SecAlign作为首个专门针对提示注入进行安全微调的基础语言模型，将攻击成功率进一步压低至4.63\%。在护栏模型方面，Inan等\cite{inan2023llamaguard}提出的LlamaGuard基于语言模型构建输入输出安全护栏，其对齐检查组件在AgentDojo上将攻击成功率从15.62\%降至1.74\%；Zeng等\cite{zeng2024shieldgemma}的ShieldGemma则基于Gemma训练通用内容审核模型，为Agent工具调用的安全对齐提供护栏支撑。然而，模型层防御存在三方面结构性局限：其一，对齐训练以单次输入-输出为粒度，难以感知跨步骤组合风险；其二，在闭源模型API占主导的产业部署环境中，Agent开发者通常无法对底层模型进行微调；其三，针对已知攻击类型微调的防御对新攻击类型的泛化能力有限，而工具调用场景的攻击面持续扩展。\textbf{第二层为架构级防御}：Debenedetti等\cite{debenedetti2025camel}提出的CaMeL采用双LLM架构将"规划"与"执行"分离，通过将不可信数据与系统指令隔离在不同模型的上下文中，首次证明架构设计可实现确定性安全保证（攻击成功率接近0\%）；Costa等\cite{costa2025fides}将经典信息流控制（IFC）理论引入Agent场景，通过对机密性与完整性标签的端到端追踪阻止不可信工具输出影响安全关键决策；Wu等\cite{wu2025isolategpt}的IsolateGPT以独立沙箱为各工具提供隔离运行环境，从操作系统层面阻断跨工具权限污染。\textbf{第三层为运行时检测与拦截}：Mou等\cite{mou2026toolsafe}提出的ToolSafe通过步骤级主动安全护栏将攻击成功率从56\%压低至1.16\%，展现了细粒度运行时守卫的有效性；An等\cite{an2025ipiguard}提出的IPIGuard则通过构建工具依赖图检测异常调用模式，为工具调用链的图结构建模提供了早期探索。\textbf{第四层为协议级与审计级防御}：Uppala\cite{uppala2026prompts}提出MCP代理以实现基于属性的架构级强制访问控制，其提出“提示指令无法可靠约束工具选择”将架构强制确立为确定性安全的唯一路径；Wang等\cite{wang2026fromagenttraces}提出Agent执行证据追踪与溯源审计的统一框架，将安全关注点从“攻击是否被阻止”延伸至“攻击是否被感知与追溯”。

目前已有多个针对Agent安全的评估基准\cite{zhang2024injectagent,balunovic2024agentdojo,zhang2025asb,li2026agentcanary}集成了直接提示注入、间接提示注入、计划投毒、记忆投毒等多种主流攻击手段，广泛地揭示了现有Agent系统在自然语言模态下指令和数据难以分离的固有模糊性导致的脆弱性，这意味着针对Agent工具调用的新威胁面应当构建新的针对性防御方法。在防御方法方面，现有防御范式主要分策略执行与系统设计、运行时检查、提示工程和微调四大类别\cite{ji2025taxonomy}，并分化为以下几类技术路线：第一类是基于信息流控制与能力约束的防御，例如CaMeL\cite{debenedetti2025camel}和FIDES\cite{costa2025fides}将可信控制流与不可信数据流分离，并通过能力机制或动态污点追踪约束敏感数据外流，从架构层面阻断注入内容对Agent规划的影响；Kim等\cite{kim2025promptflow}提出提示流完整性机制，将信息流控制理论应用于Agent提示-工具执行链，防止低权限输入通过工具调用提升权限。第二类是基于结构化输入与权限管控的防御，例如StruQ\cite{chen2025struq}和Spotlighting\cite{hines2024defending}通过结构化查询格式或标记机制将用户指令与外部数据内容在语法层面分离，降低注入内容被误解析为指令的概率；Progent\cite{shi2025progent}和SEAgent\cite{ji2026seagent}分别引入领域特定语言和强制访问控制框架实现细粒度权限管控；Zhu等\cite{zhu2025melon}提出基于掩码执行的可证明防御方法MELON，通过校验将用户指令和外部数据遮蔽前后Agent的行为一致性检测注入攻击。第三类是基于运行时监控与执行隔离的防御，例如IsolateGPT\cite{wu2025isolategpt}通过执行隔离架构将LLM推理与工具执行环境隔离，从系统架构层面降低攻击面；DRIFT\cite{li2025drift}通过动态规则生成和注入内容隔离实现运行时防御；ClawGuard\cite{zhao2026clawguard}、ToolSafe\cite{mou2026toolsafe}和ClawKeeper\cite{liu2026clawkeeper}在工具调用边界进行拦截与检查或构建多层防护体系，在AgentDojo等基准上防御成功率达86.67--94.21\%。第四类是基于语义分析与主动检测的防御，例如IPIGuard\cite{an2025ipiguard}通过构建工具依赖图建模正常调用模式，以异常依赖关系检测间接提示注入攻击；PIGuard\cite{li2025piguard}提出缓解过度防御的提示注入护栏模型，通过减少触发词偏见在检测注入攻击的同时保持低误报率；DRIFT\cite{abdelnabi2025drift}通过激活差异检测任务漂移，Attention Tracker\cite{hung2025attention}通过注意力权重分布识别注入攻击，为无需修改Agent架构的轻量级检测提供了新思路；AgentVisor\cite{ying2026agentvisor}通过语义虚拟化隔离不可信输入，AgentShield\cite{rassul2026agentshield}通过蜜罐机制在零误报率条件下检测率达90.7--100\%。然而，上述工作的评价视角大多停留于“攻击是否成功”或“防御是否阻断”，对工具调用链本身的风险结构、过程级可解释评估和事后归因机制关注不足。第五类是基于安全微调与对齐的防御，例如通过安全导向的指令微调\cite{ouyang2022training}或基于人类反馈的强化学习使模型内化安全行为边界，Wallace等\cite{wallace2024instruction}提出的指令层级训练和Chen等\cite{chen2025secalign,chen2025metasecalign}的SecAlign系列将安全对齐从通用安全扩展至工具调用场景，试图从模型参数层面抑制注入攻击的生成概率。然而，微调/对齐路线存在三方面结构性的局限：其一，对齐训练以单次输入-输出为粒度，难以感知跨步骤组合风险；其二，在闭源模型API占主导的产业部署环境中，Agent开发者通常无法对底层模型进行微调；其三，针对已知攻击类型微调的防御对新攻击类型的泛化能力有限，而工具调用场景的攻击面持续扩展。

\subsection{现有研究存在的问题}

通过上述研究现状，可以发现现有工作在以下几个方面存在不足，构成本研究的直接研究动机：

\todo{简化这里的总结，不要超过400字。}
\begin{enumerate}
  \item \textbf{G1：针对Agent工具调用链的粗放风险度量体系难以捕获过程间组合性风险，且检测与防御过程耦合不足。}
  现有针对Agent工具调用链安全问题的基准和防御方法以攻击成功率为核心指标，而对工具敏感度、权限扩散度、调用深度、不可逆操作比例、敏感数据暴露、任务必要性偏差等过程性风险指标的度量研究较少。例如，“语义权限提升”这类序列合法但组合有害的风险对现有以单次调用为粒度的访问控制系统完全不可见，也无法被任务完成率或攻击成功率等结果导向指标所量化。更重要的是，即使能够输出风险分数，现有方法也缺乏将风险等级直接映射为可操作防御动作（放行/记录/脱敏/人工确认/阻断回滚）的分级机制，导致风险感知与防御响应之间存在断层。
  
  \item \textbf{G2：Agent轨迹日志结构缺乏可信的统一标准，威胁归因与防御工作缺乏数据基础。}
  目前，Agent CoT轨迹被越来越多地用作审计日志，但其不可信性使基于CoT的审计存在根本缺陷，而防篡改的执行追踪需要从设计层面将审计日志与Agent控制流分离\cite{chu2026systematic,dehghantanha2026sokattacksurfaceagentic}。不同Agent框架记录的日志字段差异显著：部分框架只记录Prompt与Final Answer，部分记录Tool Call与Observation，鲜少完整记录计划步骤、参数来源（\texttt{src}）、数据流向（\texttt{sink}）、权限状态、人工确认标记和跨Agent委托关系\cite{chan2024visibility,microsoft2025taxonomy,wang2026fromagenttraces}，这些信息在进行审计或防御构建时由于信息不全、粒度未对齐等问题，难以使用统一的方法进行校验。Agent日志结构的不统一和不可信不仅制造了新的攻击面，也给下游依赖日志数据做安全度量、威胁归因与防御的工作形成了一定阻碍。

  \item \textbf{G3：工具调用链威胁检测局限于单步权限检查或最终结果判断，跨步骤组合风险路径识别不足，且现有行为轨迹图缺乏防御节点表示。}
  现有防御机制（如权限白名单、提示注入过滤）主要在工具调用前进行单步权限检查，或在最终输出阶段判断违规内容，但对跨步骤、跨工具的风险路径识别不足，例如“读取敏感文件—摘要压缩—调用外部发送工具”或“查询数据库—写入共享记忆—下游Agent复用”等链式风险路径\cite{dewitt2026openchallenges,debenedetti2025camel,costa2025fides}，这类链式攻击利用的是工具调用链上每次调用合法性与序列级意图之间的“语义鸿沟”，而这种差异难以通过单步行为或语义检测防御。同时，许多Agent安全基准均止步于以攻击成功率为最终指标，仅研究单次调用或有限序列检测层面，跨步骤组合风险在结构上无法被现有评估框架测量，未提供面向组合风险路径的结构化检测方法\cite{zhang2024injectagent,liu2024formalizing,chu2026systematic}。部分工作\cite{wang2025agentarmor,liu2026auditing}虽然引入了Agent多步的轨迹级审计，但输出仍为合规/违规的二元判断，无法识别或量化跨步骤组合风险路径的内部结构，仍然存在一定的局限性。此外，现有行为轨迹图方法\cite{he2025sentinelagent,wang2026aligning}均未在图结构中显式表示防御节点（$V_{\text{guard}}$）和防御边（$E_{\text{defend}}$），导致防御动作无法在图内被追踪和验证。

  \item \textbf{G4：风险发现后的威胁归因机制不足，防御动作缺乏最小干预原则，审计机制缺乏可信证据链。}
  G3指出的检测不足的问题间接导致了威胁归因和防御机制的缺失：当Agent系统出现不当工具调用或敏感数据外传时，仅输出“异常”二元标签无法满足实际审计或者防御构建的需求。因此，实际应用的LLM Agent安全防御系统必须可以解释风险路径的构成节点、关键触发工具节点位置、参数来源、风险最高的步骤，才能根据证据链定向进行防御策略的部署，而现有工具调用风险研究\cite{he2025sentinelagent,wang2026aligning,wang2026agenttrace,li2026traces}尚未提供满足上述需求的可解释威胁归因方法。更进一步，即使能够完成归因，现有防御方法也普遍采用“全量阻断”策略，缺乏在满足残余风险约束前提下最小化任务效用损失的优化机制，导致安全性与效用之间的权衡无法被系统性解决\cite{deng2025agentsunderthreat,balunovic2024agentdojo}。
  
\end{enumerate}
